\documentclass{article}
\usepackage{amsmath, amssymb, amsthm}
\usepackage{hyperref}
\usepackage{geometry}
\geometry{margin=1in}

\title{Erd\H{o}s Problem \#75}
\date{Last updated: 2025-08-31}
\author{Source: erdosproblems.com}

\begin{document}
\maketitle

\noindent\textbf{Status:} OPEN \quad \textbf{No prize}

\noindent\textbf{Tags:} graph theory, chromatic number

\medskip

\noindent\textbf{Problem Statement:}

Is there a graph of chromatic number $\aleph_1$ with $\aleph_1$ vertices such that for all $\epsilon&#62;0$ if $n$ is sufficiently large and $H$ is a subgraph on $n$ vertices then $H$ contains an independent set of size $&#62;n^{1-\epsilon}$? What about an independent set of size $\gg n$?




Conjectured by Erdős, Hajnal, and Szemerédi \cite{EHS82}. In \cite{Er95} Erd\H{o}s asks this without the condition that the graph also have $\aleph_1$ vertices, but this is an oversight, since already in \cite{EHS82} they provide such a construction. In \cite{Er95d} Erd\H{o}s offers \$1000 for a complete solution to all problems of this type (for example including also [74] ), and a 'generous reward for any significant partial results'. See also [750] .




References

[EHS82] Erd\H{o}s, P. and Hajnal, A. and Szemer\'{e}di, E., On almost bipartite large chromatic graphs . Theory and practice of combinatorics (1982), 117-123.

[Er95] Erd\H{o}s, Paul, Some of my favourite problems in number theory, combinatorics, and geometry . Resenhas (1995), 165-186.

[Er95d] Erd\H{o}s, Paul, On some problems in combinatorial set theory . Publ. Inst. Math. (Beograd) (N.S.) (1995), 61-65.


\bigskip
\noindent\small{Source: \url{https://www.erdosproblems.com/75}}
\end{document}
